% !PRINTED 653
\weberpage{653}
%<<CONT
auch die Funktion \glqq 1\grqq{} darstellbar sein muß. Nehmen wir an,
$\omega_1$ sei nicht durch $\pi$ teilbar. Nach \S\ 182, 6. gibt es Konstanten
$c_1, c_2, \ldots, c_n$, deren erste unbeschadet der Allgemeinheit $= 1$
angenommen werden kann, die den Bedingungen:
\[
\omega_1 \equiv 1, \ \omega_2 \equiv c_2, \ \ldots, \ \omega_n = c_n \quad (\mathrm{mod}\ \pi)
% [sic] printed $\omega_n = c_n$ with = instead of \equiv, the relation is a congruence modulo $\pi$ like the two preceding ones
\]
genügen. Wir setzen:
\weq{14}{
\begin{aligned}
\alpha_1 &= (z - c)\,\omega_1,\\
\alpha_2 &= \omega_2 - c_2\,\omega_1,\\
&\cdot\ \cdot\ \cdot\ \cdot\ \cdot\ \cdot\ \cdot\\
\alpha_n &= \omega_n - c_n\,\omega_1.
\end{aligned}
}

Daß dies eine Basis von $\pi$ ist, ersieht man sofort, wenn man
die Funktion (13) so darstellt:
\weq{15}{
\omega = x_2\alpha_2 + \cdots + x_n\alpha_n + (x_1 + c_2x_2 + \cdots + c_nx_n)\,\omega_1\,;
}
da $\alpha_2, \ldots, \alpha_n$ durch $\pi$ teilbar sind, so kann $\omega$ nur dann durch
$\pi$ teilbar sein, wenn die ganze Funktion in $Z$:
\[
x_1 + c_2x_2 + \cdots + c_nx_n
\]
durch $(z - c)$ teilbar und folglich $\omega$ durch die Formel (5) darstellbar
ist. Hiernach ergeben die Formeln (10) und (14):
\weq{16}{
N_a(\pi) = z - c\,,
}
also den Satz:

2. \gesperrt{Die absolute Norm eines Primfunktionals ist eine lineare Funktion von $z$.}

Dieser Satz bedeutet einen wesentlichen Unterschied zwischen
den Theorien der Zahlen und der Funktionen. In der Zahlentheorie
 gibt es Primideale verschiedener Grade, deren Norm eine
Potenz einer Primzahl ist. In der Theorie der Funktionen haben
wir nur Primfunktionale ersten Grades. Dies rührt daher, daß
wir vermöge des Fundamentalsatzes der Algebra jede ganze
Funktion einer Variablen nach dem Modul $\pi$ in lineare Faktoren
zerlegen können.

Die Zahl der Primfaktoren, in die $(z - c)$ zerfällt, ist hiernach
 $n$, da $N(z-c) = (z-c)^n$ ist.

Fassen wir die untereinander gleichen Primfaktoren in
Potenzen zusammen, so ist
\weq{17}{
z - c = \varepsilon\,\pi_1^{e_1}\pi_2^{e_2}\pi_3^{e_3}\ldots\,,
}
worin
\weq{18}{
n = e_1 + e_2 + e_3 + \cdots\,,
}
wo $\varepsilon$ eine Einheit ist.
